\section{Discussion}\label{sec:discussion}

\subsection{Why Magnitude Ranking Works}
The cell model explains why ranking by singular value matches the Fisher estimate. Each crossbar column is
quantised with its own scale (\cref{eq:quant}), and the column of component $r$ in the first crossbar holds
$\sigma_r v_r$, so its scale is $s_r=|\sigma_r|\,\|v_r\|_\infty/127$. A fixed cell noise therefore causes a weight error
proportional to $|\sigma_r|$: components with large singular values receive the largest absolute errors.
The Fisher score (\cref{eq:fisher}) multiplies this error size by the loss sensitivity of each stored value.
That it selects essentially the same components as magnitude suggests that the error size varies across
components much more than the sensitivity. The gradient score measures sensitivity to a small change of
$\sigma_r$ alone, ignoring both the error size and the errors on $u_r$ and $v_r$, and so falls behind when
errors are large.

\subsection{Relation to HyFlexPIM}
Our results agree with HyFlexPIM~\cite{song2025hyflexpim} that a hybrid SLC/MLC layout of SVD components
works and that all-MLC storage does not, but not on how to rank components. HyFlexPIM compares its
gradient ranking with magnitude- and rank-based alternatives on MRPC and CoLA at one noise level and reports
differences of about a point or less in favour of the gradient. We find no difference at mild noise
($\sigma=3\%$) and the reverse at high noise, so the findings need not conflict. Three differences in
set-up explain this:
\begin{itemize}
\item \textbf{Noise model.} Under HyFlexPIM's relative weight noise the error of a weight does not depend on
the scale of its crossbar column, so the mechanism above is absent. Its noise level is set from a bit error
rate measured on real MLC cells~\cite{fan2024mlc}, which is itself a per-cell quantity.
\item \textbf{Noise range.} HyFlexPIM evaluates a single noise level. In our sweeps the rankings differ
only at $\sigma\geq5\%$.
\item \textbf{When gradients are taken.} HyFlexPIM records gradients during its 1--3 epochs of
fine-tuning, while ours are taken on the fine-tuned model. We did not repeat their procedure, since magnitude
already matches the Fisher estimate.
\end{itemize}

\subsection{Limitations}
\begin{itemize}
\item \textbf{Scope of the error model.} Each non-ideality was varied on its own, the IR-drop model is
static, and parameter ranges were swept from harmless to failure rather than calibrated to one device, so
we report trends and comparisons rather than the accuracy of a specific chip.
\item \textbf{Known faults and cost.} Fault-aware mapping needs every chip to be tested once, and cells
that fail later need a new test and mapping. Savings are counted in memory cells and ADC conversions;
full-system energy, latency and area are not evaluated.
\item \textbf{Activation quantisation.} One 8-bit scale per tensor shifts SST-2 accuracy by up to about a
point between batch sizes and costs up to 11 points on MNLI even without errors (\cref{sec:mnli}); all
comparisons use identical batches.
\item \textbf{Training and scope.} Each training variant was run once with one training chip, as our
compute was limited to free-tier GPUs; on MNLI only placement and mapping were run. The SVD step itself costs about two points.
\end{itemize}
